设备在故障早期往往只出现肉眼难以察觉的微小振动，若未能及时发现，故障可能逐步加剧，危及人身和财产安全。压电加速度计需要逐点安装和布线，测点数量受通道限制；激光测振仪可以非接触测量，但一束激光只对准一个点，价格也很高。本作品“见微”用一台双目相机从远处观察带散斑纹理的振动表面，在 AMD Zynq-7010 上同时跟踪最多 12 个测点的亚像素位移，给出主频、峰峰值和有效值，并在现场完成声光报警、打印和移动端推送。

系统以 ALINX AX7010（XC7Z010，CLG400 封装、速度等级 $-1$）为核心。可编程逻辑（PL）承担双路 OV5640 视频采集、三路 VDMA 帧缓存、HDMI 显示、PDM 麦克风位流接收、补光 PWM 与看门狗以及打印串口；处理系统（PS）运行 Linux 与 C++17 主程序，完成 25×25 模板的逆组合 Lucas–Kanade 亚像素跟踪、静止参考点平移扣除、128 帧 Hann 窗频谱、停机基线包络门限判定、声学包络测频、界面、记录与报警；Arduino Nano ESP32 负责 NFC、SD 卡语音、Wi-Fi 与 APP 网关。我们设计了 100 mm×100 mm 的主载板、60 mm×40 mm 的前面板和 322 mm×361 mm×181 mm 的 3D 打印外壳。

2026 年 11 月 4 日的板端实测中，减速电机带动贴有散斑的铝片振动：先采集 15 s 停机基线，再测量约 10 s。相机实际帧率 45.0477 帧/s，452 帧无丢帧、无跟踪失败，帧间隔抖动 0.417 ms；12 个测点在全部 11 个分析窗中均超过 3 倍停机基线门限，逐点信号与基线之比为 36～95。视觉主频为 19.9476 Hz，麦克风仅凭声音包络测得 19.9277 Hz，两者相差 0.0199 Hz。按试验尺标换算，代表测点峰峰值 0.2090 mm、有效值 0.0701 mm，超过 0.02 mm 的报警阈值，报警灯、蜂鸣器、语音与手机 APP 同时报警。PL 工程占用 LUT 50.8\%、触发器 41.0\%、BRAM 50.8\%，已约束路径的建立时间裕量为 +0.756 ns。

\vspace{4mm}
\noindent{\sffamily\bfseries 关键词：}视觉测振；Zynq-7000；亚像素跟踪；Lucas–Kanade；停机基线；PDM 麦克风；软硬件协同
